%!TEX root = ../template.tex
%%%%%%%%%%%%%%%%%%%%%%%%%%%%%%%%%%%%%%%%%%%%%%%%%%%%%%%%%%%%%%%%%%%%
%% acronyms.tex
%% NOVA thesis document file
%%
%% Acronyms definition
%%%%%%%%%%%%%%%%%%%%%%%%%%%%%%%%%%%%%%%%%%%%%%%%%%%%%%%%%%%%%%%%%%%%

\typeout{NT FILE acronyms.tex}%


% [REV-2026-09-13][FR-01][P3][APRESENTACAO] Substituir as siglas de exemplo pelas usadas na tese
% Este ficheiro ainda contém abbrev, aaa, xpto e outras entradas do
% template. Criar apenas as necessárias: CDU, SGE, USE, EnPI, EnB, GNE,
% EII, OLS, HAC, MBB, MAE, CUSUM, PRISMA-ScR, entre outras realmente usadas.
% Definir equivalências EnPI/IDE e EnB/CER uma vez e escolher a notação
% principal. Uniformizar kt/d, t GNE/d e grafia fuel gás/fuel gas.
% [/REV-2026-09-13]
\newacronym{abbrev}{abbrev}{abbreviation of a longer text}
\newacronym{aaa}{aaa}{acronym aaa}
\newacronym{aab}{aab}{acronym aab}
\newacronym{aba}{aba}{acronym aba}
\newacronym{bbb}{bbb}{acronym bbb}
\newacronym{xpt}{xpto}{and extension of a xpto xpto xpto xpto xpto xpto xpto xpto xpto xpto xpto xpto xpto xpto xpto xpto xpto xpto xpto}

% \newacronym{novathesis}{\emph{novathesis}}{\emph{novathesis} template}
%
% \newacronym{novathesisclass}{\texttt{novathesis.cls}}{\texttt{novathesis.cls} class}


\newacronym[sort=novathesis]{novathesis}{\novathesis}{NOVAthesis \LaTeX}

\newacronym{novathesisclass}{novathesis.cls}{\novathesis\ class}

\newacronym{NOVA}{NOVA}{NOVA University Lisbon}

\newacronym{FCT}{FCT}{NOVA School of Science and Technology}

\newacronym{DI}{DI}{Department of Computer Science}
